\section{Related Work}\label{sec:relatedworks}

\BfPara{Verifier-in-the-Loop Generation} Baldur~\cite{first2023baldur}
generates and repairs whole Isabelle proofs against the verifier's own
feedback. DeepSeek-Prover~\cite{xin2024deepseek} uses a deterministic verifier
as the reward signal for large-scale synthetic proof data. Both treat the
verifier as part of the generation loop rather than as a final judge. Framing
L applies the same idea to \tlaplus{}: the verifier's diagnostic, not just its
pass/fail bit, is fed back into the next draw.

\BfPara{Grammar-Guided Decoding} Grammar-constrained decoding
restricts a language model's output to a formal grammar~\cite{geng2023grammar}.
At each step an index built from the grammar, typically a finite automaton
over the grammar's states~\cite{willard2023efficient}, masks the vocabulary
down to the tokens that keep the prefix in the language: a non-conforming
continuation is unreachable at decode time, not merely discouraged.
XGrammar~\cite{dong2024xgrammar} is the engine vLLM exposes for this, and the
dialect \gfile{} is written in. We are not aware of a
published grammar for \tlaplus{} itself. \gfile{}
targets the \tlaplus{} expression language specifically, including
constructs (nested junction lists, qualified instance application, the
symbolic operator set) that a generic programming-language grammar would not
cover. We measure it offline; we do not run it inside a generation arm.

\BfPara{\tlaplus{} Specification Synthesis} FormaLLM~\cite{bisharat2026formallm}
established the open-loop baseline this program starts from: 26.6\% SANY
parse and 8.6\% TLC model-check across 25 open-weight models. \ours{} is the
parent training program. \tlaplus{}~\cite{lamport2002specifying} and
TLC~\cite{yu1999tlc} are the target language and verifier throughout. Both
this addendum and the parent program report against industrial \tlaplus{}
use~\cite{Newcombe2015,cirstea_validating_2024} as the motivating context, not
as a benchmark.
